\chapter{Sorodna dela}
\label{sec:sorodna-dela}

To poglavje umešča nalogo med obstoječe pristope za delo s podatki sledilnika pogleda.
Sorodna dela v tem poglavju obravnavamo z vidika vprašanja, kako podatke sledilnika pogleda predstaviti za klasifikacijo čustvenih oznak in kakšno klasifikacijsko uspešnost pri tem dosegajo.
Najprej predstavimo grafovske nevronske mreže za časovne vrste in časovno-prostorske podatke, saj je osrednji del naloge pretvorba časovnega okna meritev sledilnika pogleda v graf.
Na koncu povzamemo še učenje predstavitev iz signala sledilnika pogleda ter položaj te naloge glede na sorodna dela.

\section{Grafovske nevronske mreže za časovno-\linebreak prostorske podatke}

Za modeliranje časovnih vrst se pogosto uporabljajo rekurentne nevronske mreže, kot so mreže LSTM, in transformerji, ki modelirajo časovne odvisnosti med meritvami~\cite{jin2024gnn4ts}.
Grafovska predstavitev je smiselna, kadar želimo poleg časovne urejenosti eksplicitno modelirati tudi odnose med več časovnimi vrstami ali prostorske odnose med meritvami~\cite{jin2024gnn4ts}.
Pregled GNN za časovne vrste to področje obravnava skozi naloge napovedovanja, klasifikacije, imputacije in zaznavanja anomalij.
Za to nalogo je pomembna predvsem formulacija, pri kateri posamezno časovno okno pretvorimo v grafovsko predstavitev in oznako napovemo na ravni celotnega grafa in s tem torej časovnega okna.

Pri podatkih sledilnika pogleda imajo meritve poleg časovnega vrstnega reda tudi prostorsko komponento, saj je vsaka veljavna meritev vezana na položaj pogleda na zaslonu.
Podobno ločevanje časovnih in prostorskih odvisnosti uporabljajo časovno-prostorske grafovske mreže, na primer ST-GCN pri napovedovanju prometa, kjer grafovska konvolucija modelira prostorske povezave med senzorji, časovna konvolucija pa dinamiko signala skozi čas~\cite{yu2018stgcn}.
V tej nalogi je ta pristop uporaben kot metodološka analogija, saj ne napovedujemo prihodnjih meritev, temveč klasificiramo oznako časovnega okna meritev.

\section{Učenje predstavitev iz signalov}
\label{subsec:sorodna-predstavitve}

Poleg ročno oblikovanih značilk obstajajo pristopi, ki se predstavitev naučijo neposredno iz zaporedij meritev.
S tem zmanjšajo odvisnost od ročne izbire značilk, hkrati pa so njihove predstavitve vezane na podatke in naloge, uporabljene pri predtreniranju.
V nalogi takšni modeli služijo kot most med klasičnimi modeli na ročno agregiranih značilkah in grafovskimi modeli, ki se predstavitev učijo na sami strukturi podatkov.

\ModelName{GazeMAE} je primer odprtokodnega modela, zasnovanega posebej za očesne gibe, ki z dvema časovno-konvolucijskima avtokodirnikoma ločeno kodira zaporedji koordinat pogleda in iz njih izračunanih hitrosti.
Pri samonadzorovanem učenju avtokodirnika rekonstruirata delno izpuščene vrednosti signalov~\cite{bautista2021gazemae}.
\ModelName{MOMENT} je splošnejši prednaučen temeljni model za časovne vrste, ki signal razdeli na dele fiksne dolžine in jih kodira z arhitekturo transformerjev.
Pri predtreniranju naključno zakrije nekatere dele signala in se uči njihove rekonstrukcije iz preostalega konteksta, zato je uporaben kot primerjava pri ostalih signalih~\cite{goswami2024moment}.

\section{Položaj te naloge glede na sorodna dela}

Pregled sorodnih del kaže, da se podatki sledilnika pogleda pri prepoznavanju čustvenih stanj najpogosteje uporabljajo v obliki ročno oblikovanih značilk ali zaporedij meritev~\cite{lim2020emotion,tarnowski2020eye,lu2015combining_eye_eeg}.
Grafovski modeli so pogosteje uporabljeni pri drugih vrstah časovno-prostorskih podatkov oziroma pri fizioloških signalih, predvsem pri EEG, kjer vozlišča običajno predstavljajo elektrode ali možganska področja~\cite{song2020dgcnn,zhong2022rgnn,liu2024gnn_eeg_survey}.
Obstajajo tudi pristopi, ki grafovske modele uporabljajo pri večmodalni prepoznavi čustev iz EEG in podatkov očesnih gibov~\cite{wu2021emotion,lu2015combining_eye_eeg}, vendar nismo našli dela, ki bi obravnavalo izključno podatke sledilnika pogleda kot časovno-prostorski graf posameznih meritev.

Ta naloga se zato osredotoča na primerjavo različnih predstavitev podatkov sledilnika pogleda.
Čustvene oznake uporablja kot evalvacijsko nalogo, s katero primerjamo predstavitve istega časovnega okna: agregirane značilke za klasične modele, naučene predstavitve iz prednaučenih kodirnikov~\cite{bautista2021gazemae} in časovno-prostorsko grafovsko predstavitev za GNN.

Grafovski del naloge preverja, ali eksplicitne časovne, prostorske in fiksacijske relacije med meritvami pogleda prispevajo k uporabnejši predstavitvi okna.
V predlagani predstavitvi vozlišča predstavljajo posamezne meritve, povezave pa časovno bližino, prostorsko bližino in pripadnost isti fiksaciji.
Naloga torej povezuje (1) binarno klasifikacijo čustvenih oznak iz podatkov sledilnika pogleda, (2) GNN za časovno-prostorske podatke ter (3) primerjavo ročno oblikovanih, naučenih in grafovskih predstavitev časovnih oken.
